In this section we present a self-contained derivation of the $r$-matrix structure and its main consequences, emphasizing explicit computations. The general framework follows the standard $r$-matrix formalism developed in \cite{SemenovTyanShanskii1983,Sklyanin1982,babelon2003introduction}.

A Lax pair provides a systematic way to construct conserved quantities, without making explicit use of a Poisson structure. However, Liouville integrability requires not only the existence of a Poisson structure, but also that the conserved quantities be in involution.
We now introduce a general form of the Poisson brackets between the entries of the Lax matrix which ensures this involution property.

Consider a Lax pair $L, M$ where both are $N \times N$ matrices, and assume that the matrix $L$ is diagonalizable, so that
\begin{equation} \label{eq:lambda_matrix}
    L = U \Lambda U^{-1}.
\end{equation}
The eigenvalues $\lambda_k$ of $L$ provide conserved quantities. We will not address here the question of their independence.

We introduce the following notation. Denote by $E_{ij}$ the canonical basis of the $ N \times N $ matrices, with components $(E_{ij})_{kl} = \delta_{ik} \delta_{jl}$. Then we can write
\begin{equation}
    L = \sum_{ij} L_{ij} E_{ij}.
\end{equation}
The coefficients $ L_{ij} $ of the Lax matrix are functions on the phase space. We can compute the Poisson brackets $ \{L_{ij}, L_{kl}\} $ and arrange them as follows. Define
\begin{equation}
    L_1 \equiv L \otimes 1 = \sum_{ij} L_{ij} (E_{ij} \otimes 1), 
    \quad
    L_2 \equiv 1 \otimes L = \sum_{ij} L_{ij} (1 \otimes E_{ij}).
\end{equation}
The subscript 1 or 2 indicates that the matrix $ L $ acts on the first or second factor of the tensor product, respectively. 
Likewise, for a tensor $ T $ belonging to the tensor product of two copies of $ N \times N $ matrices, we set
\begin{equation}
    T = T_{12} = \sum_{ijkl} T_{ijkl} E_{ij} \otimes E_{kl},
    \quad
    T_{21} = \sum_{ijkl} T_{ijkl} E_{kl} \otimes E_{ij}.
\end{equation}
More generally, when tensor products with several copies of $ N \times N $ matrices are involved, we denote by $ L_\alpha $ the embedding of $ L $ in the $\alpha$-th position, e.g.
\begin{equation}
    L_3 = 1 \otimes 1 \otimes L \otimes \cdots,
    \quad
    T_{\alpha\beta} \text{ the embedding of } T \text{ in positions } \alpha \text{ and } \beta.
\end{equation}
We also use $ \Tr_\alpha $ to denote the partial trace over the $\alpha$-th factor in a tensor product. For instance,
\begin{equation}
    \Tr_1 T_{12} = \sum_{ijkl} T_{ijkl} \, \Tr(E_{ij}) \, E_{kl}.
\end{equation}
\begin{definition}
    Define $ \{L_1, L_2\} $ as the matrix of Poisson brackets between the entries of $ L $:
    \begin{equation}
        \{L_1, L_2\} = \sum_{ijkl} \{L_{ij}, L_{kl}\} E_{ij} \otimes E_{kl}.
    \end{equation}
\end{definition}
For integrable systems, the Poisson brackets between the entries of the Lax matrix admit a particularly simple and structured form.

We collect a number of auxiliary results that will be used in the formulation of the $r$-matrix structure. We start with a general identity, which will serve as the basis for the subsequent lemmas and corollaries.

\begin{proposition}\label{prop:leibniz-expand-general}
    Let $L$ and $X$ be matrices whose entries belong to a Poisson algebra.
    Set $L_1 = L \otimes 1$ and $X_2 = 1 \otimes X$.
    Assume that the Poisson bracket is extended entrywise to matrix-valued functions and satisfies the Leibniz rule in each argument.
    Then, for every integer $n \ge 1$,
    \begin{equation}\label{eq:leibniz-expand-general}
        \{L_1^n, X_2\}
        =
        \sum_{p=0}^{n-1}
        L_1^{n-1-p} \{L_1, X_2\} L_1^p.
    \end{equation}
\end{proposition}

\begin{proof}
    We argue by induction on $n$.
    For $n=1$, the statement is immediate.
    Assume that \eqref{eq:leibniz-expand-general} holds for some $n$. Then, by the Leibniz rule,
    \begin{equation}
        \{L_1^{n+1}, X_2\} = \{L_1 L_1^n, X_2\} = L_1 \{L_1^n, X_2\} + \{L_1, X_2\} L_1^n.
    \end{equation}
    Using the inductive hypothesis, we obtain
    \begin{equation}
        \{L_1^{n+1}, X_2\}
        = \sum_{p=0}^{n-1} L_1^{n-p} \{L_1, X_2\} L_1^p + \{L_1, X_2\} L_1^n.
    \end{equation}
    The last term is precisely the contribution corresponding to $p=n$, hence
    \begin{equation}
        \{L_1^{n+1}, X_2\}
        = \sum_{p=0}^{n} L_1^{n-p} \{L_1, X_2\} L_1^p.
    \end{equation}
    This proves the inductive step and therefore the statement for all $n \ge 1$.
\end{proof}

\begin{corollary}\label{cor:leibniz-expand}
    Let $L$ be a matrix whose entries belong to a Poisson algebra.
    Set $L_1 = L \otimes 1$ and $L_2 = 1 \otimes L$.
    Assume that the Poisson bracket is extended entrywise to matrix-valued functions and satisfies the Leibniz rule in each argument.
    Then, for all integers $n,m \ge 1$,
    \begin{equation}\label{eq:leibniz-expand}
        \{L_1^n,L_2^m\}
        =
        \sum_{p=0}^{n-1}\sum_{q=0}^{m-1}
        L_1^{\,n-1-p}L_2^{\,m-1-q}
        \{L_1,L_2\}
        L_1^pL_2^q
    \end{equation}
\end{corollary}

\begin{proof}
    Applying Proposition \ref{prop:leibniz-expand-general} with $X_2 = L_2^m$, we obtain
    \begin{equation}
        \{L_1^n, L_2^m\}
        =
        \sum_{p=0}^{n-1}
        L_1^{n-1-p} \{L_1, L_2^m\} L_1^p.
    \end{equation}
    Expanding $\{L_1, L_2^m\}$ by the Leibniz rule in the second argument gives
    \begin{equation}
        \{L_1, L_2^m\}
        =
        \sum_{q=0}^{m-1}
        L_2^{m-1-q} \{L_1, L_2\} L_2^q.
    \end{equation}
    Substituting this into the previous formula yields
    \begin{equation}
        \{L_1^n, L_2^m\}
        =
        \sum_{p=0}^{n-1} \sum_{q=0}^{m-1}
        L_1^{n-1-p} L_2^{m-1-q}
        \, \{L_1, L_2\} \,
        L_1^p L_2^q,
    \end{equation}
    which proves the claim.
\end{proof}

\begin{proposition}\label{prop:fund}
     Let $L$ be a matrix whose entries belong to a Poisson algebra.
    Set $L_1 = L \otimes 1$ and $L_2 = 1 \otimes L$.
    Assume that the Poisson bracket is extended entrywise to matrix-valued functions and satisfies the Leibniz rule in each argument.
    Suppose that there exists an element $r$ of the tensor product algebra such that
    \begin{equation}\label{eq:fund}
        \{L_1,L_2\} = [r_{12},L_1] - [r_{21},L_2].
    \end{equation}
    Then, for all integers $n, m \ge 1$,
    \begin{equation}\label{eq:claim}
        \{L_1^n, L_2^m\} = [a_{12}^{n,m}, L_1] + [b_{12}^{n,m}, L_2],
    \end{equation}
    where
    \begin{align}
        a_{12}^{n,m} &=\sum_{p=0}^{n-1}\sum_{q=0}^{m-1} L_1^{n-1-p}L_2^{m-1-q}r_{12}L_1^{p}L_2^{q}, \\
        b_{12}^{n,m} &=-\sum_{p=0}^{n-1}\sum_{q=0}^{m-1} L_1^{n-1-p}L_2^{m-1-q}r_{21}L_1^{p}L_2^{q}.
    \end{align}
\end{proposition}
\begin{proof}
   We start from the result of the previous corollary~\ref{cor:leibniz-expand}:
    \begin{equation}
        \{ L_1^n, L_2^m \} = \sum_{p=0}^{n-1} \sum_{q=0}^{m-1} L_1^{n-1-p} L_2^{m-1-q} \{ L_1, L_2 \} L_1^{p} L_2^{q}.
    \end{equation}
    Using Corollary~\ref{cor:leibniz-expand} and substituting \eqref{eq:fund}, we obtain
    \begin{equation}
        \{ L_1^n, L_2^m \}
        =\sum_{p=0}^{n-1} \sum_{q=0}^{m-1}
        L_1^{n-1-p} L_2^{m-1-q} \left( [r_{12}, L_1] - [r_{21}, L_2] \right) L_1^{p} L_2^{q}.
    \end{equation}
    Split the sum into the two corresponding contributions. For the $r_{12}$-term we have
    \begin{align}
        \sum_{p,q} L_1^{n-1-p} L_2^{m-1-q} [r_{12}, L_1] L_1^{p} L_2^{q} &= \sum_{p,q} \left( L_1^{n-1-p} L_2^{m-1-q} r_{12} L_1^{p+1} L_2^{q} - L_1^{n-p} L_2^{m-1-q} r_{12} L_1^{p} L_2^{q} \right)\\
        &= \left[ \sum_{p,q} L_1^{n-1-p}L_2^{m-1-q} r_{12} L_1^{p}L_2^{q}\;, \;L_1 \right],
    \end{align}
    where in the second line we have used that $L_1$ and $L_2$ commute, so that the factors $L_1^{n-1-p}L_2^{m-1-q}$ and $L_1^{p}L_2^{q}$ commute with $L_1$. Hence
    \begin{equation}
        X [r_{12},L_1] Y = [X r_{12} Y, L_1],
    \end{equation}
    and we obtain
    \begin{equation}
        a_{12}^{n,m} = \sum_{p=0}^{n-1}\sum_{q=0}^{m-1}
        L_1^{n-1-p} L_2^{m-1-q} r_{12} L_1^{p} L_2^{q}.
    \end{equation}
    An analogous computation for the $-[r_{21},L_2]$ term gives
    \begin{equation}
    \sum_{p,q} L_1^{n-1-p} L_2^{m-1-q} (-[r_{21},L_2]) L_1^{p} L_2^{q}
    =
    \left[
    -\sum_{p,q} L_1^{n-1-p} L_2^{m-1-q} r_{21} L_1^{p} L_2^{q},
    L_2
    \right].
\end{equation}
Combining the two commutators yields \eqref{eq:claim}.


    
\end{proof}

\begin{lemma} \label{lemma:TrL^n involution}
    If (\ref{eq:claim}) holds then the quantities $\Tr(L^n)$ are in involution, i.e.
    \begin{equation}
        \{\Tr(L^n), \Tr(L^m)\} = 0 \quad \forall \, n, m.
    \end{equation}
\end{lemma}
\begin{proof}
    Take the trace over both spaces:
    \begin{equation}
        \{\Tr\left(L^n\right) , \Tr\left(L^m\right)\}
        = \Tr_{12}\left(\left\{L_1^n , L_2^m\right\}\right).
    \end{equation}
    Using the linearity 
    \begin{equation}
        \{\Tr\left(L^n\right) , \Tr\left(L^m\right)\} 
        = \Tr_{12}\left(\left[a_{12}^{n,m} , L_1\right]\right) + \Tr_{12}\left(\left[b_{12}^{n,m} , L_2\right]\right).
    \end{equation}
    By cyclicity of the trace
    \begin{equation}
        \Tr_{12}\left(\left[X , Y\right]\right) = 0,
    \end{equation}
    hence
    \begin{equation}
        \{\Tr\left(L^n\right) , \Tr\left(L^m\right)\} = 0,
    \end{equation}
    this shows that the spectral invariants $\Tr(L^n)$ are in involution.
\end{proof}


\begin{corollary} \label{cor:importante}
    Assume there exists an element $r_{12}$ of the tensor product algebra such that
    \begin{equation}
        \{L_1, L_2\} = [r_{12}, L_1] - [r_{21}, L_2].
    \end{equation}
    Then the spectral invariants $\Tr (L^n)$ are in involution, i.e.
    \begin{equation}
        \{\Tr(L^n), \Tr(L^m)\} = 0 \quad \, \forall \, n,m.
    \end{equation}
\end{corollary}

\begin{proof}
    This follows directly from Proposition~\ref{prop:fund} together with Lemma~\ref{lemma:TrL^n involution}.
\end{proof}

We now address the converse problem: starting from the involution of the eigenvalues of $L$, we show that, locally on phase space, one can reconstruct an $r$-matrix structure of the form (\ref{eq:r_matrix}).

\begin{proposition}[Local reconstruction of an $r$-matrix]\label{prop:existence_rmatrix}
    Assume that $L$ is locally diagonalizable as in (\ref{eq:lambda_matrix}) (for instance, with simple spectrum), and that its eigenvalues Poisson commute,
    \begin{equation}
        \{\lambda_i,\lambda_j\}=0 \quad \forall \, i,j.
    \end{equation}
    Then, locally on phase space, one can reconstruct an element $r_{12}$ of the tensor product algebra such that
    \begin{equation}\label{eq:r_matrix}
        \{L_1, L_2\} = [r_{12}, L_1] - [r_{21}, L_2].
    \end{equation}
\end{proposition}

\begin{proof}[Sketch of proof]
We outline the construction locally on a region of phase space where $L$ is diagonalizable with simple spectrum, so that one can choose $U$ smoothly such that $L = U \Lambda U^{-1}$.

Using the Leibniz rule, we expand
\begin{equation}
    \{L_1,L_2\} = \{U_1 \Lambda_1 U_1^{-1}, U_2 \Lambda_2 U_2^{-1}\}.
\end{equation}
All terms can be expressed in terms of $U$, $\Lambda$, and the brackets $\{U_1,U_2\}$, $\{U_1,\Lambda_2\}$, $\{\Lambda_1,U_2\}$, using the identity
\begin{equation}
    \{U^{-1},X\} = - U^{-1} \{U,X\} U^{-1}.
\end{equation}

Conjugating the whole expression by $U_1^{-1} \otimes U_2^{-1}$, one reduces to computing the Poisson brackets of $\Lambda$ together with the transformed brackets of $U$. The terms involving $\{\Lambda_1,\Lambda_2\}$ vanish by assumption, since the eigenvalues Poisson commute.

The remaining contributions can be organized into commutators with $\Lambda_1$ and $\Lambda_2$, yielding an expression of the form
\begin{equation}
    \{L_1,L_2\} = [r_{12},L_1] - [r_{21},L_2],
\end{equation}
where $r_{12}$ is constructed locally from $U$ and its Poisson brackets.

The construction is therefore valid on any open set where the diagonalization of $L$ can be chosen smoothly.
\end{proof}
